\section{问题重述与分析}

\subsection{问题背景}

顺流和逆流的声传播时间并不相同，超声波流量计正是利用这一差异求取轴向流速。对流速作时间与截面积分，就可以得到窗口体积。单声道在理想剖面下能够给出有效估计，实际管路却很少保持在理想状态。弯头、阀门、缩径段或其他扰流件会改变速度分布，偏流与旋流也随之出现。

已有工作从不同角度处理这一误差。Mousavi 等借助半三维仿真求取几何剖面修正因子；Sakhavi 和 Nouri 用广义速度剖面降低积分结果对流场偏离的敏感程度。Papathanasiou 等关注声道布置与扰流修正的配合，Zhao 等则用 ANN 融合多路径观测\cite{Mousavi2020FPCF,Sakhavi2022GeneralizedProfile,Papathanasiou2022FlowDC,Zhao2014ANN}。这些研究为本文提供了一个清晰起点：物理积分负责确定流量主尺度，剖面特征与受约束补偿用于处理偏差。

本题给出一台五声道超声波流量计的窗口化测试数据，包含标准体积、基础积分估计量、五条等效声道观测量、A/B 方向差异、剖面特征、启停动态特征和零点状态特征。建模目标是在不使用标准体积、日期、窗口编号等非运行时信息作为在线输入的前提下，输出每个窗口的估计体积，并在无扰流线性度、重复性和扰流综合指标上达标。

\subsection{数据概况}

主数据表记录了 $159$ 个窗口，每个窗口包含 $40$ 个字段。D0 无扰流参考样本只有 $10$ 个；其余 $149$ 个窗口来自 D1 至 D8 八种扰流条件。数据覆盖的名义流量点为 20、30、40、50、60、70 和 80。题目评价脚本计算部分指标时采用 40 至 100 的有效范围。

窗口内原始观测序列共有 $29322$ 行，能够通过窗口编号回接主表。另一个几何参数表给出了等效半径、截面积、五条声道的位置与映射关系，同时提供多套固定积分权重。

泛化评价采用三层验证。首先，以日期、流量点和扰流编号构成分层单元，形成 $118$ 个训练窗口与 $41$ 个验证窗口的标签盲示例划分；其次，固定模型配置，在 $40$ 个不同随机种子的同类分层划分上统计性能波动；最后，将完整分层单元留作验证，实施整组留出压力测试。划分过程只利用样本的分组信息，不查看标准体积。模型应用于新日期时，应预先固定特征、核参数和校准形式，并使用独立数据检验。

\subsection{问题拆解}

问题一要求从传播时间差模型推导单声道体积估计公式，并解释多声道积分权重的物理含义；问题二要求在无扰流数据上建立多声道标定模型，并通过相对误差、组均值误差和组内重复性评价其稳定性；问题三要求分析扰流下的剖面、A/B 差异、动态和零点特征，构造不依赖人工扰流编号的在线识别与补偿方法；问题四要求整合前三问，形成固定规则和参数的最终窗口体积估计模型，并用题目给定评价脚本输出窗口误差、组达标数和三项综合指标。
